\subsection{Control implementation and scaling}
\label{sec:control-implementation}
The live prototype and the policy simulations have distinct performance costs.
The implementation revisions described here address live control overhead; they
do not revise the reported simulation outcomes or establish hardware throughput.

\paragraph{Incremental event consumption.}
Previously, live JSONL draining reused the full-file offline reader, so each tick
parsed and reapplied historical records. A per-consumer byte cursor now separates
incremental consumption from full replay. With bounded event size, parsing costs
$O(\Delta E)$ for newly appended events rather than $O(E)$ for the entire history;
an unchanged log needs only a metadata check. A new board still replays from the
beginning to reconstruct state. Complete malformed lines are counted and skipped,
partial last lines are retried, and observed truncation or file replacement resets
the cursor. Rotation requires coordination to avoid losing an unread old-file tail;
unobserved copy-truncate/regrowth is unsupported. This cursor does not checkpoint
model state or acknowledge successful downstream application. Crash-safe recovery
still requires consistent state/log checkpoints and duplicate handling. Consumed
arrival records are discarded while future timestamps are retained, preserving
the existing exclusion of starts older than the previous board tick.

The design prioritizes event semantics before CPU kernels. Tests exercise repeated
drains and board ticks, partial and malformed lines, independent readers, file
replacement, observed truncation, and I/O retry. Component benchmarks and their
source hashes accompany the development documentation; none measures GPU service
capacity or end-to-end tail latency.

\paragraph{Batched control delivery.}
Previously each hold caused a serial HTTP request and a fresh queue scan. The
controller now sends bounded batches of at most 1,024 holds. The proxy validates
all records before applying a batch, discards expired updates, and performs one
admission pass. For $D$ directives, backlog $Q$, and batch bound $B$, the repeated
list-scan term changes from $O(DQ)$ to $O(D+\lceil D/B\rceil Q)$, excluding release
selection. Round trips decrease from $D$ to $\lceil D/B\rceil$; payload work stays
$O(D)$. Chunks are independent, not a durable whole-plan transaction. Updates
retain their original meaning: they govern subsequent submissions and do not
retime already queued requests. Versioned deltas and automatic retry are deferred
because expiry renewal and restart synchronization require a separate protocol.
The per-request board client also retrieves service and next-tool estimates in
one request, with local fallback on unavailable results. CPU work still shares
the board event loop; these changes do not establish request-tail isolation.

\paragraph{Indexed admission.}
The live queue previously scanned its pending list for every individual completion.
It now maintains priority and FCFS ready heaps, delayed-release and promotion
heaps, and request-identity dictionaries. Stable insertion sequences resolve ties;
versioned records invalidate old priorities, and periodic linear-time rebuilds
bound stale heap storage. Ordinary release costs $O(\log Q)$ amortized, reducing
individual backlog drainage from $O(Q^2)$ to $O(Q\log Q)$; simultaneous due events
still require processing each affected request. Overflow transitions retain the
original hold-dropping and FCFS policy and can cost $O(Q)$. A single wake-up timer
tracks holds and promotions. Randomized differential tests compare release order
and state with the former scan implementation. Diagnostic counts and peer-rank
queries still scan the backlog, so this is not a logarithmic bound on the whole
HTTP request path. The simulator's worker queue is a separate implementation.
